%------------------------------------------------------------------------------%
% GENERATORS OF STRONGLY CONTINUOUS SEMIGROUPS                                 %
%------------------------------------------------------------------------------%
% File  : Hille_Yosida.tex                                                     %
% Author: Klaus Hoeflinger, Beethovenstrasse 11, D-73117 Wangen                %
% Email : khoeflinger@t-online.de                                              %
%------------------------------------------------------------------------------%

%------------------------------------------------------------------------------%
%                       DOCUMENT CLASS and INCLUDE FILES                       %
%------------------------------------------------------------------------------%
\documentclass[11pt,twoside]{article}
\usepackage{geometry}
\geometry{paperheight=241mm,
          paperwidth=152mm,
          top=22mm,
          bottom=17mm,
          left=20mm,
          right=20mm,
          textwidth=112mm,
          textheight=202mm,
          headheight=10mm,
          headsep=3mm,
          footskip=9mm,
          includehead,
          includefoot,
          marginparsep=0mm,
          marginparwidth=0mm}

\renewcommand{\baselinestretch}{0.945}\normalsize

\AtBeginDocument{\setlength\abovedisplayskip{2mm}}
\AtBeginDocument{\setlength\belowdisplayskip{2mm}}

%------------------------------------------------------------------------------%
% define title formats                                                         %
%------------------------------------------------------------------------------%
\usepackage{titlesec}

\renewcommand{\thesection}{\arabic{section}}
\renewcommand{\sfdefault}{FiraSans}
\renewcommand{\ttdefault}{pcr}
\renewcommand{\rmdefault}{antiqua}

\titleformat{\part}[display]
{\normalfont \large \rmfamily \bfseries \centering}
{\small{\thepart.}}{2pt}{}

\titleformat{\section}[runin]
{\normalfont \rmfamily \bfseries}
{\thesection}{1pt}{.\;}[.]

%\titleformat{\section}
%{\normalfont \bfseries \small \filcenter}
%{\thesection.\;}{0mm}{}

\titleformat{\subsection}[runin]
{\normalfont \itshape}
{}{0mm}{}

\titleformat{\subsubsection}[runin]
{\normalfont \itshape}
{}{}{}{}

\titleformat{\paragraph}[runin]
{\normalfont}
{}{}{}{}

\titleformat{\subparagraph}[runin]
{\normalfont}
{}{}{}{}

\titlespacing*{\part}          { 0pt}{ 0pt}{ 8pt}
\titlespacing*{\section}       { 0pt}{ 7pt}{ 4pt}
\titlespacing*{\subsection}    { 0pt}{14pt}{ 6pt}
\titlespacing*{\subsubsection} { 0pt}{ 6pt}{12pt}
\titlespacing*{\paragraph}     { 0pt}{ 6pt}{12pt}
\titlespacing*{\subparagraph}  { 0pt}{ 0pt}{12pt}

%------------------------------------------------------------------------------%
% define enumerate parameters                                                  %
%------------------------------------------------------------------------------%
\usepackage{enumitem}
\setlength{\parindent}{3pt}
\setlength{\parskip}  {0pt}

%\setlist{noitemsep}

\setlist[enumerate]{
         leftmargin=12pt,
         rightmargin=0pt,
         itemsep=1pt,
         itemindent=3pt,
         topsep=3pt,
         labelsep=4pt,
         partopsep=0pt,
         labelwidth=0pt,
         listparindent=0pt,
         parsep=0pt}

\setlist[itemize]{
         leftmargin=12pt,
         rightmargin=0pt,
         itemsep=0pt,
         itemindent=1pt,
         topsep=1pt,
         labelsep=4pt,
         partopsep=1pt,
         labelwidth=0pt,
         listparindent=0pt,
         parsep=0pt}

%------------------------------------------------------------------------------%
% define packages                                                              %
%------------------------------------------------------------------------------%
\usepackage{upref}
\usepackage{amssymb}
\usepackage{slantsc}
\usepackage{amsmath}
\usepackage{amsthm}
\usepackage{thmtools}
\usepackage{amscd}
\usepackage{verbatim}
\usepackage{latexsym}
\usepackage{multicol}
\usepackage{graphicx}
\usepackage{setspace}
\usepackage[ngerman,english]{babel}
\usepackage[protrusion=true,expansion=true]{microtype}
\usepackage{tikz}
\usetikzlibrary{arrows,arrows.meta,decorations.markings}
\usepackage{mathrsfs}
\usepackage{amsfonts}
\usepackage{datetime2}
\usepackage{textgreek}
\usepackage{hyperref}
\usepackage{pifont}
\usepackage{relsize}
%\usepackage{biblatex}
\relsize{-0.05}
\usepackage[skip=0mm]{parskip}
\usetikzlibrary{decorations.pathmorphing}

\usepackage{mlmodern}
\usepackage[T5,T1]{fontenc}
\usepackage{ragged2e}

%\usepackage{mathabx}
%\usepackage{times}

%\usepackage{fontspec}
%\setmainfont{Nimbus Roman No9 L}
%\usepackage[T1]{fontenc}

%------------------------------------------------------------------------------%
% define page layout                                                           %
%------------------------------------------------------------------------------%
\usepackage{fancyhdr}
\pagestyle{fancy}
\setlength{\headheight}{30.0pt}
\renewcommand{\headrulewidth}{0pt}

\AtBeginDocument{
  \addtolength\abovedisplayskip{-0.0\baselineskip}
  \addtolength\belowdisplayskip{-0.0\baselineskip}
}

\fancyhead{}
\fancyfoot{}
\addtolength{\headwidth}{0.02\marginparwidth}

\makeatletter
\let\ps@plain\ps@fancy
\makeatother

\renewcommand\sectionmark[1]
{
 \def\sectionname{#1}
 \markright{\thesection #1}
}

\makeatletter
\@addtoreset{section}{part}
\makeatother

%------------------------------------------------------------------------------%
% define footnote without number                                               %
%------------------------------------------------------------------------------%
\newcommand{\myfootnote}[1]{%
\renewcommand{\thefootnote}{}%
\footnotetext{#1}%
\renewcommand{\thefootnote}{\arabic{footnote}}%
}

%------------------------------------------------------------------------------%
% define numbering in equations                                                %
%------------------------------------------------------------------------------%
\numberwithin{equation}{section}

%------------------------------------------------------------------------------%
% define newtheorem                                                            %
%------------------------------------------------------------------------------%
\declaretheoremstyle[
headfont=\scshape, 
headindent = 0mm, %\parindent,
%postheadhook = {\hspace{0mm}\newline},
%postheadspace = \newline, 
spaceabove = 3mm, 
spacebelow = 3mm,
numberwithin=section,
name=Theorem]{khMATH}
\declaretheorem[style=khMATH]{theorem}

\declaretheoremstyle[
headfont=\bfseries, 
headindent = 0mm, %\parindent,
%postheadhook = {\hspace{0mm}\newline},
%postheadspace = \newline, 
spaceabove = 3mm, 
spacebelow = 3mm,
numberwithin=part,
name=Theorem]{khMATH1}
\declaretheorem[style=khMATH1]{theorem1}

\declaretheoremstyle[
headfont=\scshape, 
headindent = 0mm, %\parindent,
%postheadhook = {\hspace{0mm}\newline},
%postheadspace = \newline, 
spaceabove = 3mm, 
spacebelow = 3mm,
numberwithin=section,
name=Corollary]{khMATH1}
\declaretheorem[style=khMATH1]{corollary}

%            name         counter     label              numeration-mode
\theoremstyle{plain}
\newtheorem*{introduction}            {\sc Introduction}
\newtheorem {standard}                {}                 [section]
\newtheorem*{definition}              {\sc Definition}
\newtheorem{satz1}                    {\sc Satz}
\newtheorem{lemma}                    {\sc Lemma}
\newtheorem*{proposition}             {\sc Proposition}
%\newtheorem{corollary}               {\sc Corollary}
\newtheorem*{corollar}                {\sc Korollar}
\newtheorem*{remark}                  {\sc Remark}
\newtheorem*{remarks}                 {\sc Remarks}
\newtheorem*{example}                 {Example}
\newtheorem*{examples}                {Examples}
\newtheorem*{theo}                    {\sc Theorem}

%------------------------------------------------------------------------------%
% define tabular                                                               %
%------------------------------------------------------------------------------%
\usepackage{tabularx}
\newcolumntype{L}[1]{>{\raggedright\arraybackslash}p{#1}}
\newcolumntype{C}[1]{>{\centering\arraybackslash}  p{#1}} 
\newcolumntype{R}[1]{>{\raggedleft\arraybackslash} p{#1}} 

%------------------------------------------------------------------------------%
% define \sh, \zB                                                              %
%------------------------------------------------------------------------------%
\def\dsh{{d.\,h.}}
\def\ie{{i.\,e.}}
\def\zB{z.\,B.}
\renewcommand{\abstractname}{ }

%------------------------------------------------------------------------------%
% change default spacing for cases                                             %
%------------------------------------------------------------------------------%
\makeatletter
\renewcommand*\env@cases[1][1.6]{%
  \let\@ifnextchar\new@ifnextchar
  \left\lbrace
  \def\arraystretch{#1}%
  \array{@{}l@{\quad}l@{}}%
}
\makeatother

\newenvironment{rcases}
  {\left.\begin{aligned}}
  {\end{aligned}\right\rbrace}

%------------------------------------------------------------------------------%
% change default for \predisplaypenalty                                        %
%------------------------------------------------------------------------------%
\predisplaypenalty=990
\allowdisplaybreaks \numberwithin{equation}{section}

%------------------------------------------------------------------------------%
% general notations                                                            %
%------------------------------------------------------------------------------%
\def\*{\star}
\def\={\,=\,}
\def\+{\,+\,}
\def\xsubset{{\,\subset\,}}
\def\xtimes{{\,\times\,}}
\def\xeof{{\,\in\,}}
\def\eq{{\,=\,}}
\def\xto{{\,\to\,}}
\def\eqdef{\,:=\,}
\def\:{{\,:\,}}
\def\ele{{\,\in\,}}
\def\exists{\mbox{ exists }}
\def\and{\mbox{ and }}
\def\for{\mbox{ for }}
\def\fa{\mbox{ for all }}
\def\fe{\mbox{ for each }}
\def\qfa{\quad \mbox{ for all }}
\def\df{\,\dotfill}
\def\iff{if and only if }
\def\wrt{with respect to}

\makeatletter
\newcommand \Df {\leavevmode \cleaders \hb@xt@ .70em{\hss .\hss }\hfill \kern \z@}
\makeatother

\def\sql{\mathfrak{a}}               % sesquilinear form a
\def\DF{B}                           % Dirichlet form a
%------------------------------------------------------------------------------%
% number sets and special elements                                             %
%------------------------------------------------------------------------------%
\def\P{\mathbb{H}}                   % half plane of $\C$
\def\J{I}                            % interval I
\def\N{{\mathbb{N}}}                 % unsigned integers
\def\Q{{\mathbb{Q}}}                 % rational integers
\def\Z{{\mathbb{Z}}}                 % integers
\def\R{{\mathbb{R}}}                 % real numbers
\def\Rn{{\mathbb{R}}^n}              % real numbers in Rn
\def\Rd{{\mathbb{R}}^d}              % real numbers in Rn
\def\C{{\mathbb{C}}}                 % complex numbers
\def\F{{\mathbb{K}}}                 % arbitrary field
\def\Torus{{\mathbb{T}}}             % complex circle line
\def\T{\tau}                         % upper bound of interval (0,t)
\def\sign#1{\operatorname{sign}(#1)} % sign of a number

%------------------------------------------------------------------------------%
% set theory                                                                   %
%------------------------------------------------------------------------------%
\def\G{G}                            % domain $G$
\def\clG{\overline{\G}}              % closure of $G$
\def\bndG{\partial \G}               % boundary of $G$
\def\setA{\mathcal{A}}               % set A
\def\setB{\mathcal{B}}               % set B
\def\setC{\mathcal{C}}               % set C
\def\setM{M}                         % set M
\def\setN{N}                         % set N
\def\setS{\mathcal{S}}               % set S
\def\famF{\mathfrak{F}}              % family of sets
\def\indI{\mathcal{I}}               % index set I
\def\pot#1{\mathfrak{P}(#1)}         % power set

\def\relR{\mathbf{R}}                % relation R
\def\mapf{f}                         % mapping f
\def\mapg{g}                         % mapping g

\def\set#1#2{\{ \, #1 \,: \, #2 \, \}}
\def\tensor#1#2{#1 \otimes #2}
\def\Tens{\oplus}

\def\cal#1{{\mathcal{#1}}}
\def\aut#1{\operatorname{Aut}(#1)}

\def\cross#1#2{#1 {\,\times\,} #2}
\def\multicross#1#2{#1 \times  \ldots \times  #2}
\def\multitens#1#2 {#1 \otimes \ldots \otimes #2}

%------------------------------------------------------------------------------%
% group theory                                                                 %
%------------------------------------------------------------------------------%
\def\1{{\mathsf{1}}}                % unit element of a monoid
\def\0{{\mathsf{0}}}                % unit element of an Abelian monoid
\def\kernel#1{\mathfrak{N}(#1)}     % kernel of a homomorphism
\def\cokernel#1{\mathfrak{C}{(#1)}} % cokernel of a homomorphism
\def\Hom#1{\operatorname{Hom}(#1)}

%------------------------------------------------------------------------------%
% linear algebra                                                               %
%------------------------------------------------------------------------------%
\def\vecV{\mathit{V}}               % vector space V
\def\vecH{\mathit{H}}               % vector space H
\def\vecW{\mathit{W}}               % vector space W
\def\vecU{\mathit{U}}               % subspace U
\def\convC{\mathcal{C}}             % convex set C
\def\basH{\mathfrak{B}}             % Hamel base B

\def\LO#1{\mathscr{L}(#1) }         % algebra of linear transformations

\def\scalar#1#2{ \langle #1,#2 \rangle }
\def\trace#1{\operatorname{tr}(#1)}
\def\dim#1{{\operatorname{dim} \, #1 }}
\def\codim#1{{\operatorname{codim} \, #1 }}
\def\dim{\operatorname{dim}}
\def\codim{\operatorname{codim}}
\def\ind{\operatorname{ind}}

\def\genG{\cal{G}}

\def\algA{\cal{A}}
\def\algB{\cal{B}}
\def\algU{\cal{U}}

\def\idI{\cal{I}}

%------------------------------------------------------------------------------%
% topology                                                                     %
%------------------------------------------------------------------------------%
\def\compK{{K}}          % compact set C
\def\openO{\mathcal{O}}  % open set O
\def\openU{\mathcal{U}}  % open set U
\def\U{\mathcal{U}}      % neighbourhood U
\def\V{\mathcal{V}}      % neighbourhood V
\def\d{\mathit{d}}       % metric function d

\def\interior#1{\mathring{#1}}
\def\closure#1{{\overline{#1}}}

\def\topT{\mathcal{T}}   % topology T
\def\topX{X}             % topological space X
\def\topY{Y}             % topological space Y
\def\topZ{Z}             % topological space Z
\def\metrS{M}            % metric space M

\def\grpG{G}             % topological group G
\def\grpA{A}             % topological group A
\def\grpB{B}             % topological group B
\def\grpC{C}             % topological group C
\def\grpH{H}             % topological group H
\def\grpU{U}             % topological group U
\def\grpV{V}             % topological group V
\def\topG{G}             % topological group G
\def\topH{H}             % topological group H
\def\topU{U}             % topological subgroup U

\def\subS{\cal{S}}
\def\riR{R}              % ring R

%------------------------------------------------------------------------------%
% calculus                                                                     %
%------------------------------------------------------------------------------%
\def\supp#1{\operatorname{supp}(#1)}
\def\singsupp#1{\operatorname{sing \, supp}(#1)}

\def\re{\operatorname{Re}}
\def\im{\operatorname{Im}}
\def\grad{\operatorname{grad}}

%\def\abs#1 {\left|#1\right|}
\def\abs#1 {\lvert #1 \rvert}
%\def\abs#1 {\left|\,#1\,\right|}

%------------------------------------------------------------------------------%
% measure theory                                                               %
%------------------------------------------------------------------------------%
\def\salgA{\Sigma}               % sigma-algebra S
\def\salgB{\cal{B}}              % sigma-algebra B
\def\Borel#1{\cal{B}_{#1}}       % Borel-algebra 
\def\LB{\lambda}                 % Lebesgue-Borel measure

%------------------------------------------------------------------------------%
% function spaces                                                              %
%------------------------------------------------------------------------------%
\def\Lp{L^p(\G)}                 % spaces L_p
\def\Lq{L^q(\G)}                 % spaces L_q
\def\Lh{L^2(\G)}                 % spaces L_2
\def\Li{L^{\infty}(\G)}          % spaces L_infty
\def\Lc{L_{loc}^1(\G)}           % spaces L_1^{loc}
\def\Ck{C^k(\G)}                 % spaces C_k

\def\BV#1{BV(#1) }               % set of functions of bounded variation
\def\AC#1{AC(#1) }               % set of absolutely continuous functions

\def\MP#1{\mathscr{M}^+(#1) }    % set of positive measures
\def\MS#1{\mathscr{M}^-(#1) }    % vector space of finite signed measures
\def\MC#1{\mathscr{M}(#1) }      % vector space of complex measures
\def\MCR#1{\mathscr{M}_{r}(#1) } % vector space of regular complex measures

\def\SobW{W^{k,p}(\G)}                % Sobolev space W
\def\SobO{\mathring{W}^{k,p}(\G)}     % Sobolev space W
%\def\WN#1#2{\mathring{W}^{#1}(#2) }   % Sobolev space W
\def\WN#1#2{W_0^{#1}(#2) }   % Sobolev space W
%\def\HN#1#2{\mathring{H}^{#1}(#2) }   % Sobolev space H
\def\HN#1#2{H_0^{#1}(#2) }   % Sobolev space H
%------------------------------------------------------------------------------%
% functional analysis                                                          %
%------------------------------------------------------------------------------%
\def\snorm#1{\lVert #1 \rVert}
\newcommand{\norm}[1]{\left\lVert #1 \right\rVert}

\def\topV{V}                       % topological vector space V
\def\topW{W}                       % topological vector space W
\def\normS{X}                      % normed space

\def\banX{\mathit{X}}              % Banach space X
\def\banY{\mathit{Y}}              % Banach space Y
\def\banZ{\mathit{Z}}              % Banach space Z
\def\dbanX{\mathit{X^*}}           % dual space of Banach space X
\def\dbanY{\mathit{Y^*}}           % dual space of Banach space Y
\def\dbanZ{\mathit{Z^*}}           % dual space of Banach space Z

\def\banA{\mathit{A}}              % Banach algebra A
\def\banB{\mathit{B}}              % Banach algebra B
\def\banC{\mathit{C}}              % C*-algebra C


\def\domain#1{\mathfrak{D}(#1)}    % domain
\def\range#1{\mathfrak{R}(#1)}     % range

\def\isoJ{\mathfrak{J}}            % isometry J
\def\repA{{\cal{A}_{\sql}}}        % representation operator A
\def\linS{{S}}                     % linear operator S
\def\linG{{G}}                     % linear operator G
\def\linL{{L}}                     % linear operator L
\def\linT{{T}}                     % linear operator A
\def\linA{\mathit{A}}              % linear operator A
\def\linB{{B}}                     % linear operator B
\def\linM{{M}}                     % linear operator M
\def\linU{{U}}                     % linear operator U
\def\linI{{I}}                     % linear operator I
\def\A{\cal{A}}                    % linear differential operator A
\def\BDO{\cal{B}}                  % linear boundary differential operator B
\def\linU{{U}}                     % unitary operator U
\def\linP{{P}}                     % projection P
\def\compA{k}                      % compact operator k
\def\linFR{f}                      % finite rank operator f
\def\fredA{A}                      % Fredholm operator F

\def\HAM{\cal{H}}                  % Hamilton operator
\def\PA{\partial^{\alpha}}         % elementary differential operator
\def\DO{\cal{A}}                   % linear differential operator
\def\BC{\cal{B}}                   % boundary operator
\def\SLO{\cal{L}_{p,q}}            % Sturm-Liouville operator
\def\MO{\cal{L}_{M}}               % momentum operator

\def\HS#1{S(#1)           }        % algebra of Hilbert-Schmidt operators
\def\FR#1{F(#1)           }        % algebra of finite rank operators
\def\TC#1{N(#1)           }        % algebra of trace class operators
\def\BU#1{\mathscr{U}(#1) }        % algebra of unitary linear operators
\def\BO#1{\mathscr{L}(#1) }        % algebra of bounded linear operators
\def\CO#1{\mathscr{C}(#1) }        % algebra of compact linear operators
\def\CL#1{\mathscr{C}(#1) }        % algebra of compact linear operators
\def\FO#1{\Phi(#1) }               % set of Fredholm operators
\def\FK#1#2{\Phi_{#2}(#1) }        % set of Fredholm operators with index k

\def\hilH{\mathit{H}}              % Hilbert space H
\def\clU{\mathit{U}}               % closed subspace U

\def\orthO{\mathfrak{O}}           % orthonormal system O
\def\orthS{\mathfrak{S}}           % orthonormal system S
\def\basB{\mathfrak{B}}            % basis B of a Hilbert space

\def\GL#1 {\mathbf{GL}(#1)}        % linear group
\def\OG#1 {\mathbf{O}(#1)}         % linear group
\def\SOG#1 {\mathbf{SO}(#1)}       % linear group

%------------------------------------------------------------------------------%
% distribution theory                                                          %
%------------------------------------------------------------------------------%
\def\disF{F}                  % distribution F
\def\disG{G}                  % distribution G
\def\disH{H}                  % distribution H
\def\disU{U}                  % distribution U
\def\disT{T}                  % distribution T
\def\SF{C^{\infty}(\G)}       % space of smooth functions
\def\TF{C_c^{\infty}(\G)}     % space of compactly supported smooth functions
\def\TFRd{C_c^{\infty}(\R^d)} % space of test functions
\def\SRd{\mathscr{S}(\R^d)}   % Schwartz space
\def\SR1{\mathscr{S}(\R)}     % Schwartz space
\def\B1#1{C^1(#1) }           % space of continuously differentiable functions
\def\Bm#1{C^m(#1) }           % space of continuously differentiable functions
\def\Ck{C^k(\G)}
\def\TD{\mathscr{S}'(\R^d)}   % space of temperated distributions
\def\E{\cal{E}}               % space of distributions with compact support
\def\D{\mathscr{D}}           % D
\def\FT{\mathcal{F}}          % Fourier transformation F
\def\FT{\mathscr{F}}          % Fourier transformation F

%------------------------------------------------------------------------------%
% other definitions                                                            %
%------------------------------------------------------------------------------%
\def\Proof{\textsc{Proof}}
\def\FRECHET{Fr\'{e}chet}
\def\ARZELA{Arzel\'{a}}
\def\GARDING{G\r{a}rding}
\def\HOELDER{H\"older}
\def\SCHROEDINGER{Schr\"odinger}
\def\JOERGENS{J\"orgens}
\def\WUEST{W\"uest}
\def\KAEHLER{K\"aehler}
\def\HOEFLINGER{H\"oflinger}
\def\POINCARE{Poincar\'{e}}
\def\FA{Functional Analysis}

\def\Author   {K.\,{\HOEFLINGER}}
\def\AuthorUC {K.\,HOEFLINGER}
\def\Version  {1.0}
\def\Email    {khoeflinger@t-online.de}


%\renewcommand\thesection{\Roman{section}}

\def\Title{Generators of Strongly Continuous Semigroups}

\titleformat{\section}
{\normalfont \bfseries \filcenter}
{\thesection.\;}{0mm}{}

\titleformat{\subsection}%[runin]
{\normalfont \itshape}
{}{}{}

\titlespacing*{\section}   {0pt}{12pt}{4pt}
\titlespacing*{\subsection}{0pt}{6pt}{0pt}

%------------------------------------------------------------------------------%
%                                BEGIN OF DOCUMENT                             %
%------------------------------------------------------------------------------%
\begin{document}
\thispagestyle{empty}

%\renewcommand{\ttdefault}{pcr}

\centerline{{\textit{Topics in Spectral and Semigroup Theory}}}
\vspace{22pt}
\centerline{\large{\textbf{\Title}}}
\vspace{4pt}
\centerline{\small{{\textsc{\Author}}}}
\vspace{15pt}

\fancyhead[C] {\small{\textsc{\Title}}}
\fancyhead[OL]{\small{\thepage}}
\fancyhead[ER]{\small{\thepage}}

\myfootnote
{
  Version {\Version} from \today;
  Email:  \textit{khoeflinger@\,t-online.de}
}

%------------------------------------------------------------------------------%
%                                 INTRODUCTION                                 %
%------------------------------------------------------------------------------%

\def\SG{$\{S(t)\}_{t \geq 0}$}
\def\TG{$\{T(t)\}_{t \geq 0}$}
\def\GG{$\{G(t)\}_{t \geq 0}$}
\def\EG{$\{\cal{E}(t)\}_{t \geq 0}$}

\section{Introduction}
%------------------------------------------------------------------------------%
In the forties of the 20$^{\textrm{th}}$ century
the cornerstones of semigroup theory has been developed independently
by the mathematicians \textsc{E.\,Hille} and \textsc{K.\,Yosida}.
Both proved in different ways a theorem that clarifies the condition,
when a linear operator $\linA$ acting in a Banach space $\banX$
is the infinitesimal generator of a strongly continuous semigroup {\SG}
of bounded linear operators.

\subparagraph{}
This subject is of great importance {\wrt} the existence and uniqueness
of solutions for the abstract initial-value problem
\[
\dot{x} + \linA x = \0, \quad x(0) = x_0 \in \banX,
\]
which is meaningful for evolutional differential equations
posed in a domain $\G \subset \Rd$
and governed by linear-elliptic partial differential operators.

\fancyhead[OC]{\small{\thesection.} \textsc{\sectionname}}
\fancyhead[EC]{\small{\textsc{\Title}}}

\section{Strongly Continuous Semigroups}
%------------------------------------------------------------------------------%
Let $\banX$ be a complex Banach space and $\BO{\banX}$ the unital Banach algebra
of bounded linear operators $\linS \: \banX \xto \banX$
with identity $\1_{\banX}$.
Then by a \textit{C$_0$-semigroup} or
a \textit{strongly continuous semigroup} of bounded linear operators
{\SG} is meant an operator-valued mapping
\[
\linS \: [0,\infty) \xto \BO{\banX},
\]
shortly denoted {\SG},
which distinguishes by the conditions

\begin{itemize}[leftmargin=12mm]
\item[(G-1)]
$\linS(t+s) \eq \linS(t) \circ \linS(s)$ for all $t,s \geq 0$
(the \textit{semigroup property});

\item[(G-2)]
$\lim_{\,t \downarrow 0} \, \linS(t)\,x \eq x$ for all $x \xeof X$.
\end{itemize}

\paragraph{}
It follows from (G-2),
which is nothing but strong continuity from the right in $t \eq 0$,
that $\linS(0) \eq \1_{\banX}$.
The conditions (G-1) and (G-2) together
imply continuity from the right of the semigroup even in each $t > 0$.
In case of $\banX \eq \R$,
both properties are shared by the family {\EG} of exponential functions
$\cal{E}(t) \: \lambda \xeof \R \mapsto e^{-\lambda t}$,
hence {\EG} serves as an elementary example of a strongly continuous semigroup.

\subparagraph{} %- classification of strongly continuous semigroups
For a given Banach space $\banX$
the totality of strongly continuous semigroups may be classified as follows:
for each such {\SG} there are constants $M \geq 0$ and $\omega \xeof \R$
such that
\begin{equation}\label{X_SEMIGROUP_GROWTH}
\norm{\linS(t)} \, \leq \, M e^{\omega t}
\quad \for t > 0.
\end{equation}

In fact,
let $M \geq 1$ such that $\norm{S(s)} \leq M$ holds for all $s \xeof [0,1]$.
For $t > 0$ there is $s \xeof [0,1)$ and $n \xeof N$ such that $t \eq n{+}s$.
Then
\[
\norm{\linS(n)} \leq M \norm{S(n-1)}
\]
for $n > 1$ and
\[
\norm{\linS(t)}
\begin{comment}
\= \norm{\linS(n{+}s)} \= \norm{\linS(n) \circ \linS(s)}
\]
\[
\end{comment}
\leq \norm{\linS(n)} \cdot \norm{\linS(s)} \leq
M^{n{+}1} \= M e^{n \cdot \ln{M}}.
\]

\subparagraph{} %- semigroups of contractions
In the special case
$(M,\omega) \eq (0,1)$ and $\norm{S(t)} \leq 1$ for $t \geq 0$,
respectively,
the semigroup is said to be \textit{contracting}.

\paragraph{} %- orbit maps
By the \textit{orbit map} or the \textit{trajectory} $\eta_x$
of a C$_0$-semigroup {\SG}
associated with an initial element $x \xeof \banX$
we mean the vector-valued mapping $\eta_x \: [\,0,\infty) \xto \banX$
defined by $\eta_x(t) := S(t) \, x$.
Both conditions (G-1) and (G-2) together imply
that $\eta_x$ is continuous.
But the orbit map is even differentiable under a very mild assumption.
In fact,
$\eta_x$ is differentiable in $(0,\infty)$
for each $x \xeof \banX$
iff
it is differentiable from the right in $t \eq 0$ only.

\section{Infinitesimal Generators}
%------------------------------------------------------------------------------%
If {\SG} is a strongly continuous semigroup,
then its \textit{generating operator} or \textit{infinitesimal generator}
(briefly denoted the \textit{generator})
is defined as a mapping $\linA$ in $\banX$ by means of the graph
\[
\Gamma_{\linA}
\eqdef
\{(x,y) \in X \xtimes X \:
y \= \lim_{\,h \downarrow 0} \, \frac{1}{h}(\linS(h)x-x) \mbox{ exists}\}.
\]
Hence the \textit{domain of definition} $\domain{\linA}$ of the mapping $\linA$
consists of those elements $x \xeof \banX$,
for which the orbit map $\eta_x$ is differentiable from the right in $t \eq 0$,
or equivalently,
\[
\domain{\linA} \eqdef
\{x \xeof \banX \: \eta_x\:[0,\infty) \xto \banX \mbox{ is differentiable}\}.
\]
It is not hard to show that
$\domain{\linA}$ is a linear subspace of $\banX$
and the mapping $\linA$ is linear,
so $\linA$ constitutes a "linear operator" in $\banX$.

\paragraph{} %- theorem 3.1
The following theorem contains a couple of assertions,
which will later be used to prove the main theorems of this text.

\begin{theorem}\label{XXX}
Let {\SG} be a C$_0$-semigroup with infinitesimal generator $\linA$.
Then
\begin{itemize}[leftmargin=30pt]

\item[(i)]
For $x \xeof \banX$
\begin{equation}\label{XXX_1}
x_t \eqdef \lim_{h \to 0} \frac{1}{h} \int_t^{t+h} S(s)x \,ds \= S(t)x.
\end{equation}

\item[(ii)]
For $x \xeof \banX$ and $\int_0^t S(s)x\,ds \xeof \domain{\linA}$
\begin{equation}\label{XXX_2}
\linA \int_0^t S(s)x \, ds \= S(t) - x.
\end{equation}

\item[(iii)]
For $x \xeof \domain{\linA}$ and $S(t)x \xeof \domain{\linA}$
\begin{equation}\label{XXX_3}
\frac{d}{dt} S(t) x \= \linA \,S(t)x \= S(t)\, \linA x.
\end{equation}

\item[(iv)]
For $x \xeof \domain{\linA}$
\begin{equation}\label{XXX_4}
S(t)x - S(s)x \=
\int_s^t S(\tau)\,\linA x\,d\tau \=
\int_s^t \linA\,S(\tau)x \, d\tau.
\end{equation}
\end{itemize}
\end{theorem}

\subparagraph{}
\Proof:
Assertion (i) follows from the continuity of the orbit mapping $\eta_x$.
To show assertion (ii), let $x \xeof \banX$ and $h > 0$. We have
\[
\frac{S(h)-\1_X}{h} \int_0^t S(s)x\,ds \=
\frac{1}{h} \int_0^t  S((s+h)-S(s)x)\,ds \=
\]
\[
\frac{1}{h} \int_t^{t+h}  S(s)x \, ds - \frac{1}{h} \int_0^h S(s)x \,ds.
\]
Now, taking $h \downarrow 0$,
the right hand side tends to $S(x)-x$ whence (ii) follows.\\
To prove (iii), let $x \xeof \domain{\linA}$ and $h > 0$. Then
\begin{equation}\label{XXX_31}
\frac{S(h)-\1_X}{h} S(t) x \=
S(t) \frac{S(h)-\1_X}{h} x \to S(t) \,\linA x \mbox{ as } h \downarrow 0.
\end{equation}
Hence $S(t)x \xeof \domain{\linA}$ and $\linA\,S(t)x \eq S(t)\linA x$.
From \eqref{XXX_31} we also deduce that
\[
\frac{d^+}{dt} S(t) x \= \linA S(t)x \= S(t) \,\linA x.
\]
Now we need to verify that for $t>0$ the left derivative exists and
equals to $S(t)\,\linA x$.
To do so, we write
\[
\lim_{\downarrow 0}\,(\frac{S(t)x-S(t-h)x}{h} - S(t)\,\linA x) \=
\]
\[
\lim_{\downarrow 0}\,S(t-h)\, (\frac{S(h)x-x}{h} - \linA x) \, + \,
\lim_{\downarrow 0}\,(S(t-h)\,\linA x -S(t) \, \linA x)
\]
and notice that both terms of the last line are zero.
The first one due to the fact that $x \xeof \domain{\linA}$ and
$\norm{S(t-h)}$ is bounded on $0 \leq h \leq t$,
and the second one because of property (G-2) of $S(t)$,
which implies assertion (iii).\\
Finally,
assertion (iv) is deduced by integrating \eqref{XXX_3} from $0$ to $t$.
\qed

\paragraph{} %- closed linear operators
Let us turn to the concept of densely defined and closed linear operator.
First,
a linear operator $\linA$ is said to be \textit{densely defined} in $\banX$,
if its domain of definition $\domain{\linA}$ lies dense in $\banX$.
Second,
the operator $\linA$ is called \textit{closed},
if the set
\[
\Gamma_{\linA} \eqdef
\{\,(x,\linA x) \:x \xeof \domain{\linA}\,\} \; \xsubset \; \banX \xtimes \banY,
\]
denoted the \textit{graph} of $\linA$,
forms a closed subset of the space $\banX \xtimes \banY$
{\wrt} the product topology $\topT_{\,\banX {\times} \banY}$.
This property especially implies that the \textit{kernel}
$\kernel{\linA} := \{x \xeof \domain{\linA}: \linA x \eq \0\}$
is closed.

\subparagraph{} %- sequentially closed operators
It is a matter of fact
that a linear operator $\linA$ acting between normed linear spaces is closed
iff it is \textit{sequentially closed},
that is,
$x_n \xto x$ and $\linA x_n \xto y$ implies
$x \xeof \domain{\linA}$ and $\linA x \eq y$,
briefly written
\[
\begin{rcases}
 x_n       \xto x \quad \\
 \linA x_n \xto y \quad \\
\end{rcases}
\quad \Longrightarrow \quad
x \xeof \domain{\linA} \, \and \, \linA x \eq y.
\]
This can be verified by the fact
that a subset of a normed linear space is closed
\iff it is sequentially closed.
Checking the sequentially closeness is the usual method
to prove that a given linear operator is graph-closed.

\subparagraph{}
We finally mention that if $\linA$ is closed,
then for $\zeta \xeof \C$ the operator $\zeta-\linA$
with domain of definition $\domain{\zeta-\linA} \eq \domain{\linA}$ is closed,
too.
Moreover,
if $\linA$ is closed and bijective,
then the inverse operator
$\linA^{-1}\: \range{\linA} \xsubset \banX \xto \domain{\linA} \xsubset \banX$
is also linear and closed.
The assertion of $\linA^{-1}$ being closed is based on the facts
that\\
(1) if $\Gamma_A$ is a closed set in $\banX \xtimes \banY$,
then $\Gamma_{A^{-1}}$ is closed in $\banY \xtimes \banX$ and\\
(2) the permutation $(x,y) \mapsto (y,x)$ is a homeomorphism
between the spaces $\banX \xtimes \banY$ and $\banY \xtimes \banX$.

\paragraph{} %- infinitesimal generators are densely defined and closed
Our next theorem indicates that infinitesimal generators fit into
the class of densely defined and closed linear operators:

\begin{theorem}\label{GENERATORS_ARE_DENSELY_DEFINED_AND_CLOSED}
If $\linA$ is the generator of a strongly continuous semigroup,
then $\linA$ is densely defined and closed.
\end{theorem}

\Proof:
To show that $\linA$ is densely defined,
consider the value $x_t$ as defined in \eqref{XXX_1}.
From assertion \ref{XXX}\,(ii) we deduce $x_t \xeof \domain{\linA}$
for all $t > 0$ and by \ref{XXX}\,(i) we have $x_t \to x$ as $t \downarrow 0$,
which implies that the closure of $\domain{\linA}$ is the whole of $\banX$.

\subparagraph{}
To show that $\linA$ is closed,
let $x_n \xeof \domain{\linA}$ with $x_n \xto x$ and $\linA x_n \xto y$.
By \ref{XXX}\,(iv) we have
\begin{equation}\label{GDDC_1}
S(t)x_n -x_n \= \int_0^t S(t)\,\linA x_n \, ds,
\end{equation}
where the integrand converges to $S(s)y$ uniformly on bounded intervals.
Thus letting $n \to \infty$ in \eqref{GDDC_1}, we obtain
\[
S(t)x \,-\, x \= \int_0^t S(t)y \, ds.
\]
Dividing the last identity by $t>0$ and making $t \downarrow 0$,
we deduce by employing \ref{XXX}\,(i) that $x \xeof \domain{\linA}$
and $\linA x \eq y$.
\qed

\paragraph{} %- the uniqueness theorem for generators
Using the fact that infinitesimal generators are densely defined,
we may quote the following uniqueness theorem:

\begin{theorem}
If $\linA$ and $\linB$ are generators of C$_0$-semigroups,
then $\linA \eq \linB$ imply that the semigroups generated by them
are identical.
\end{theorem}

\subparagraph{}
\Proof:
Let {\TG} and {\SG} be the generated semigroups of $\linA$ and $\linB$,
respectively,
and let $x \xeof \domain{\linA} \eq \domain{\linB}$.\\
By \ref{XXX}\,(iii),
the function $s \mapsto T(t-s)\,S(s)x$ is differentiable and
\[
\frac{d}{ds} T(t-s)\, S(s)x \= -\linA\, T(t-s)\,S(s)x + T(t-s)\,\linB\, S(s)x
\]
\[
\= -T(t-s)\,\linA \,S(s)x + T(t-s)\,\linB\,S(s)x \= \0
\]
holds.
Therefore the function $s \mapsto T(t-s)\,S(s)x$ 
is constant and in particular its values at $s \eq 0$
and $s \eq t$ are the same,
that is,
$T(t)x \eq S(t)x$.
This is true for any $x \xeof \domain{\linA}$,
since by theorem \ref{GENERATORS_ARE_DENSELY_DEFINED_AND_CLOSED}
the domain of definition $\domain{\linA}$ is dense in $\banX$
and $T(t)$ and $S(t)$ are bounded,
thus $T(t)x \eq S(t)x$ for every $x \xeof \banX$.
\qed

\section{C$_0$-Semigroups and Cauchy Problems}
%------------------------------------------------------------------------------%
The \textit{homogeneous Cauchy problem} of first order associated with
a densely defined and closed linear operator $\linA$ acting in a Banach space
$\banX$ is to find a solution of the initial value problem 

\begin{equation}\label{HOMOGENEOUS_CAUCHY_PROBLEM}
\begin{cases}[1.2]
\quad \dot{x}(t) \= \linA \, x(t),\\
\quad x(0) \= x_0 \xeof \banX.\\
\end{cases}
\end{equation}
The task is to find \textit{vector-valued} solutions $x \: [0,\T) \xto \banX$
of $\dot{x} \eq \linA x$ with $x(0) \eq x_0$.

\paragraph{} %- solution concepts for the Cauchy problem
Remember that a vector-valued function $x \: [a,b] \xto \banX$
is called \textit{classical differentiable} in $t_0 \xeof (a,b)$,
if there is an element $x'(t_0) \xeof \banX$ such that
\[
\lim_{\,t \xto t_0} \,
\norm{ \, \frac{x(t)-x(t_0)}{t-t_0}\, - x'(t_0) \, } \= 0.
\]
Moreover,
the function $x$ is called {classical differentiable in} $(a,b)$,
if $x$ is so in each $t \xeof (a,b)$.
The linear space consisting of functions having \textit{continuous}
classical derivatives in $(a,b)$ is denoted $C^1(a,b;\banX)$.

\subparagraph{} %- spaces of k-times differentiable functions
The spaces $C^k(a,b;\banX)$ ($k \xeof \N$) are defined recursively,
and the space $C^{\infty}(a,b;\banX)$ is defined
by taking the intersection of all spaces $C^k(a,b;\banX)$,
$k \xeof \N_0$.

\paragraph{} %- solution concepts for the Cauchy problem
Concepts of solution for problem \eqref{HOMOGENEOUS_CAUCHY_PROBLEM}
are given as follows:

\begin{itemize}[leftmargin=12pt]
\item[1.]
A vector-valued function
$x \xeof C([0,\T),\domain{\linA}) \cap C^1((0,\T),\banX)$
is called a \textit{classical solution} of \eqref{HOMOGENEOUS_CAUCHY_PROBLEM},
if it satisfies the initial condition $x(0) \eq x_0$
and the differential equation in each point $t \xeof (0,\T)$;

\item[2.]
a continuous function $x \: [0,\T) \xto \banX$ is said to be
a \textit{mild solution} of \eqref{HOMOGENEOUS_CAUCHY_PROBLEM} if
\[
\mbox{(1) } \, \tilde{x}(t) \eqdef
\int_{\,0}^t x(s) \, ds \, \xeof \domain{\linA}
\quad \and \quad
\mbox{(2) } \, \linA \, \tilde{x}(t) \eq x(t)-x_0
\]
holds for all $0 {\,<\,} t {\,<\,} \T$,
where the integral on the right-hand side of (1) 
is meant in the Riemannian sense.
One is led to the notion of mild solution in a natural way
by integrating both sides of the equation $\dot{x} \eq \linA \, x$,
which then becomes into
\[
\int_{\,0}^t \dot{x} ds \=
x(t)-x(0) \= 
\int_{\,0}^t \linA \, x(s) \, ds.
\]
If $\linA$ is closed,
then the action of $\linA$ and performing the integration commutes
for continuous functions $x \: [0,t] \xto \domain{\linA}$,
that is,
\[
\int_{\,0}^t \linA \, x(s) \, ds  \= \linA \, \int_{\,0}^t x(s) \, ds,
\]
and to the end we obtain the equation
\[
x(t) \= x(0) + \linA \, \int_{\,0}^t x(s) \,ds
\quad \quad (t \xeof (0,\T)).
\]
\end{itemize}

\subparagraph{} %- comparing classical and mild solutions
Notice that each classical solution is a mild solution,
and reversely a mild solution is classical
iff it is continuously differentiable in $(0,\T)$.
Thus classical and mild solutions of homogeneous Cauchy problems
only differ {\wrt} this regularity property.

\paragraph{} %-- correctly and well-posed Cauchy problems
After having the notions of classical and mild solution for the homogeneous
Cauchy problem available,
we want to introduce the concepts of \textit{correct posedness} and
\textit{well-posedness}.

\subparagraph{} %- definitions of correctly and well-posed problems
First,
the homogeneous Cauchy problem \eqref{HOMOGENEOUS_CAUCHY_PROBLEM}
is called \textit{correctly posed},
if there is a unique classical solution $x$
for any initial value $x_0 \xeof \domain{\linA}$.
Second,
it is denoted \textit{well-posed},
if it is correctly posed and
the solution $x$ depends continuously on the initial value $x_0$,
or equivalently,
if for each $\T {\,>\,} 0$ there is $c_{\T} {\,>\,} 0$ such that
\[
\norm{x(t)} \leq c_{\T} \norm{x_0}
\]
for all $t \xeof (0,\T)$ and all $x_0 \xeof \banX$.

\subparagraph{}
In \cite{EN}
the continuous dependency of the solutions from the initial values
is expressed as follows:
if $\{x_n\}_{n \xeof \N}$ is a converging sequence in $\domain{\linA}$
with limit element $\0 \xeof \banX$,
then the related solutions $t \mapsto x_n(t)$ are uniformly convergent
to the zero function in all compact intervals $[0,t_0]$
for $0 {\,<\,} t_0 < \T$.

\subparagraph{}
We mention without proof
that problem \eqref{HOMOGENEOUS_CAUCHY_PROBLEM} is well-posed
iff
it is correctly posed and the operator $\linA$ is \textit{regular},
that is,
$\rho(\linA) \ne \emptyset$ holds
(for the definition of the set $\rho(\linA)$ see section 5).

\paragraph{}
The main theorem of this section provides the following equivalence
and emphasizes the meaning of C$_0$-semigroups for the solution theory
of abstract linear evolution equations:

\begin{theorem}
The Cauchy problem \eqref{HOMOGENEOUS_CAUCHY_PROBLEM} is well-posed
iff
the linear operator $\linA$ is the infinitesimal generator of a C$_0$-semigroup.
\end{theorem}

\subparagraph{}
For the proof of this theorem see \cite{EN},
theorem II.6.6 at p.\,112.


\section{Resolvent Functions}
%------------------------------------------------------------------------------%
In this section we want to introduce some required concepts of spectral theory.
Any complex number $\zeta \xeof \C$ is called a \textit{regular value}
of the operator $\linA$,
if $\linA - \zeta \: \domain{\linA} \xto \banX$ is bijective
and the inverse operator
\[
R(\linA,\zeta) \eqdef (\zeta -\linA)^{-1} \: \banX \xto \domain{\linA}
\]
is continuous.
The set $\rho(\linA)$ of regular values is denoted
the \textit{resolvent set} of $\linA$.
If $\zeta \xeof \rho(\linA)$,
then $R(\linA,\zeta)$ is everywhere defined in $\banX$.
Since we assume $\linA$ to be closed,
the operators $\zeta -\linA $ and $R(\linA,\zeta)$ are closed,
too,
and the \textit{closed graph theorem}
implies that $R(\linA,\zeta)$ is continuous,
that is,
$R(\linA,\zeta) \xeof \BO{\banX}$\footnote{
\,If $\linA$ is \textit{not} closed,
one has in addition to require that
for $\zeta \xeof \rho(\linA)$ the inverse operator $R(\linA,\zeta)$
is bounded.}.
The linear operator $R(\linA,\zeta)$ is called the \textit{resolvent}
of $\linA$ {\wrt} $\zeta \xeof \rho(\linA)$,
and the operator-valued mapping
\begin{equation}\label{VI_RESOLVENT_FUNCTION}
\begin{cases}[1.2]
\quad R(\linA,\cdot) \: \rho(\linA) \xto \BO{\banX}\,,\\
\quad \zeta \mapsto (\zeta - \linA)^{-1} \\
\end{cases}
\end{equation}
is denoted the \textit{resolvent function} (or shortly the \textit{resolvent})
of $\linA$.

\subparagraph{}
But the resolvents of a linear operator $\linA$ are not independent
from each other.
In particular,
for all $\zeta_1,\zeta_2 \xeof \rho(\linA)$
the \textit{first resolvent equation}
\begin{equation}\label{VI_RESOLVENT_EQUATION}
R(\linA,\zeta_1) - R(\linA,\zeta_2)
\=
(\zeta_2-\zeta_1) \, R(\linA,\zeta_1) \circ R(\linA,\zeta_2)
\end{equation}
holds,
which implies that
the resolvents $R(\linA,\zeta_1)$ and $R(\linA,\zeta_2)$
commute with each other.

\paragraph{}
Let us turn to the resolvent function of a linear operator $\linA$
with non-empty resolvent set $\rho(\linA)$.
In fact,
if $\zeta_0 \xeof \rho(\linA)$,
we have $\zeta \xeof \rho(\linA)$ for all $\zeta \xeof \C$ satisfying
\[
\abs{\zeta - \zeta_0 } \, < \, r \eqdef \norm{R(\linA,\zeta_0)}^{-1},
\]
thus $\rho(\linA)$ is an open set in $\C$.
For values $\zeta$ in the disk with center $\zeta_0$ and radius $r$
one may represent the resolvent function $R(\linA,\cdot)$ as a power series
with coefficients in the operator algebra $\BO{\banX}$:
\[
R(\linA,\zeta) \=
\sum_{n=0}^{\infty} \, R(\linA,\zeta_0)^{n+1} \, (\zeta_0-\zeta)^n.
\]
This implies that
$R(\linA,\cdot)$ is locally holomorphic in $\rho(\linA)$.

\paragraph{}
Our next theorem ephasizes the meaning of operators to be assumed as closed:

\begin{theorem}
If $\rho(\linA)\,{\neq}\,\emptyset$, then $\linA$ is closed.
\end{theorem}

\subparagraph{}
\Proof:
Let $\zeta \xeof \rho(\linA)$ such that
$\linA-\zeta$ is bijective with domain of definition
$\domain{(\linA-\zeta)^{-1}} \eq \banY$ and range
$\range{(\linA-\zeta)^{-1}} \eq \domain{\linA}$.
Let $x_n \xto x$ with $x_n \xeof \domain{\linA}$ for $n \xeof \N$
and $\linA x_n \xto y$.
Then
\begin{equation}\label{XXX_TOOL}
x \eq \lim_{n \to \infty} x_n \eq
\lim_{n \to \infty} (\zeta-\linA)^{-1}(\zeta-\linA) x_n \eq
(\zeta-\linA)^{-1}(\zeta x-y).
\end{equation}
Thus $x \xeof \range{(\zeta-\linA)^{-1}} \eq \domain{\linA}$.
When applying $(\zeta-\linA)^{-1}$ from left to \eqref{XXX_TOOL},
we obtain $\zeta x -\linA x \eq \zeta x -y$,
hence $\linA x \eq y$,
which is nothing but closeness of $\linA$.
\qed

\paragraph{}
Thus a reasonable spectral theory makes only sense
for linear operators being at least continuous or closed.
Those operators,
which are closed and in addition own non-empty resolvent sets,
are often called \textit{regular}.

\paragraph{} %- Laplace transforms
In the second part of this section we want to establish a relationship between
strongly continuous semigroups and the resolvents of their infinitesimal operators,
denoted the \textit{Laplace transform} of the semigroup.

\subparagraph{}
Let the semigroup {\SG} be strongly continuous.
Then for fixed $x \xeof \banX$ and $\zeta \xeof \C$,
the value $L_n(\zeta,x)$ defined by
\begin{equation}\label{LAPLACE}
L_n(\zeta,x) \eqdef
\frac{1}{(n-1)!} \int_{\,0}^{\infty} t^{n-1}e^{-\zeta t}S(t)x \,dt
\end{equation}
is called the (right side) \textit{Laplace transform} of order $n$
of the semigroup {\wrt} $x$ and $\zeta$.
The expressions have to be understood as
improper vector-valued Riemann integrals.
But we cannot expect that $L_n(\zeta,x)$ exists without further assumptions.

\paragraph{}
The following theorem describes necessary conditions for a linear operator
$\linA$ to be the infinitesimal generator of a C$_0$-semigroup:

\begin{theorem}
Let {\SG} be a strongly continuous semigroup of type $(M,\omega)$
and with infinitesimal generator $\linA$. Then:

\begin{itemize}[leftmargin=15pt]
\item[(a)]
If $\re \zeta > \omega$,
then $\zeta \xeof \rho(\linA)$ and $R(\linA,\zeta)^nx \eq L_n(\zeta,x)$
for all $n \xeof \N$ and $x \xeof \banX$.

\item[(b)]
For all $n \xeof \N$ and $\re \zeta > \omega$
the following resolvent estimates hold:
\[
\norm{R(\linA,\zeta)^n} \, \leq \, M (\re \zeta - \omega)^{-n}.
\]
\end{itemize}
\end{theorem}

\subparagraph{}
\Proof:
The values $L_n(\zeta,x)$ are well-defined for $\re \zeta > \omega$
and $x \xeof \banX$.
Hence the mapping $L_n(\zeta)x := L_n(\zeta,y)$ forms a linear operator
in $\banX$ with norm estimate
\[
\norm{L_n(\zeta,\cdot)}
\, \leq \,
\frac{1}{(n-1)!} \int_0^{\infty} t^{n-1} e^{-\re \zeta t}e^{\omega t}\,dt
\= \frac{M}{(\re \zeta - \omega)^n}.
\]
holds.
For $n \eq 1$ we have $L_1(\zeta)\linA x \eq -x + \zeta L_1(\zeta) x$,
whereas for $n > 1$ the equality
$L_n(\zeta)\linA x \eq -L_{n-1}(\zeta)x + \zeta L_1(\zeta)x$ holds.

Applying $(S(h)-\1_X)/h$ to equation \eqref{LAPLACE},
we obtain by taking the limit $h \to 0$\,:
\[
\begin{cases}[1.2]
\quad \linA\,L_1(\zeta)x \= -x + \zeta L_1(\zeta)x\\
\quad \linA\,L_n(\zeta)x \= -L_{n-1}(\zeta)x + \zeta L_n(\zeta)x \quad (n>1),\\
\end{cases}
\]
which leads to
\[
\begin{cases}[1.2]
\quad L_1(\zeta)(\zeta-\linA) \= (\zeta-\linA)L_1(\zeta) \= \1_X,\\
\quad L_n(\zeta)(\zeta-\linA) \= (\zeta-\linA)L_n(\zeta) \= L_{n-1}(\zeta).\\
\end{cases}
\]
In conclusion,
the resolvent $R(\linA,\zeta)$ exists and for all $n \xeof \N$
we have $R(\linA,\zeta)^n \eq L_n(\zeta)$.
Especially for $\re \zeta > \omega$
we have $(\omega,\infty) \xsubset \rho(\linA)$.
\qed

\paragraph{}
Let us put together the necessary conditions for the linear operator $\linA$
to be an infinitesimal generator:

\begin{corollary}
In order $\linA$ to be the generator of a C$_0$-semigroup of type $(M,\omega)$,
the following properties must be fulfilled.

\begin{itemize}[leftmargin=15pt]
\item[(a)]
$\linA$ is densely defined and closed.
\item[(b)]
$(\omega,\infty) \xsubset \rho(\linA)$ and
\begin{equation}\label{RESOLVENT_ESTIMATION_HY}
\norm{R(\linA,\zeta)^n} \leq M (\zeta - \omega)^{-n}
\quad \for \zeta \in (\omega,\infty) \and n \xeof \N.
\end{equation}
\end{itemize}
\end{corollary}

\subparagraph{}
The family of inequalities quoted in (b) of this corollary is often denoted
the \textit{Hille-Yosida condition}.

\subparagraph{}
In the case, where $(M,\omega) \eq (1,\omega)$,
that is,
the semigroup is \textit{quasi-contractive},
we have $\norm{R(\linA,\zeta)} \leq (\zeta-\omega)^{-1}$ for $\zeta > \omega$ ($n \eq 1$),
which implies that all the conditions for $n > 1$ are also fulfilled.
Hence the family of conditions in \eqref{RESOLVENT_ESTIMATION_HY} reduces to
a single one, and this can often verified very easily.
Due to the importance of this fact,
we quote this special case in a further corollary:

\begin{corollary}
In order $\linA$ to be the generator of a C$_0$-semigroup of type $(1,\omega)$,
the following properties must be fulfilled.

\begin{itemize}[leftmargin=15pt]
\item[(a)]
$\linA$ is densely defined and closed.
\item[(b)]
$(\omega,\infty) \xsubset \rho(\linA)$ and
\begin{equation}\label{RESOLVENT_ESTIMATION_HY2}
\norm{R(\linA,\zeta)} \leq \frac{1}{\zeta - \omega}
\quad \for \zeta \in (\omega,\infty).
\end{equation}
\end{itemize}
\end{corollary}

\section{Generator Theorems of Hille-Yosida Type}
%------------------------------------------------------------------------------%
As it has already been elaborated in the third section,
to each strongly continuous semigroup {\SG} of bounded operators
there is a linear operator $\linA$ as its infinitesimal generator.
Reversely,
the question arises,
under which conditions a linear operator constitutes
the infinitesimal generator of a C$_0$-semigroup.
For bounded operators the answer can be given quickly:
if $\norm{\linA}$ is the operator norm of $\linA$,
then the mappings
\[
\exp(\linA)(t)\,x \eqdef \sum_{n=0}^{\infty} \frac{(t \linA)^n}{n!} \,x
\]
for $x \xeof \banX$ and $t>0$ form a strongly continuous semigroup.

\paragraph{} %- Hille-Yosida theorem
The general answer,
however,
is given by the \textit{Feller-Miyadera-Hille-Yosida} \textit{theorem}:

\begin{theorem}\label{GENERAL_HY-THEOREM}
Let $\linA$ be a linear operator acting in a complex Banach space $\banX$
with domain of definition $\domain{\linA}$.
Then the following assertions are equivalent to each other:

\begin{itemize}[leftmargin=15pt]
\item[(a)]
The operator $\linA$ generates a C$_0$-semigroup of type $(M,\omega)$.

\item[(b)]
the operator $\linA$ is densely defined and closed,
the resolvent set of $\linA$ fulfills $(\omega,\infty) \xsubset \rho(\linA)$
and for $\zeta > \omega$ the resolvent $R(\linA,\zeta)$
satisfies the Hille-Yosida conditions

\begin{equation}\label{HILLE-YOSIDA_CONDITIONS}
\norm{R(\linA,\zeta)^n}
\, \leq \, M (\zeta - \omega)^{-n} \quad (n \in \N).
\end{equation}
\end{itemize}
\end{theorem}

\subparagraph{}
\Proof:
It was already shown in the last section that (a) implies (b)
(see corollary 5.1).
Reversely,
it is to prove that (b) implies (a).
First of all,
we shall reformulate and prove theorem \ref{GENERAL_HY-THEOREM}
for the case of quasi-contractive semigroups
and then come back to the proof of the implication (b) $\Rightarrow$ (a)
at the end of this section.

\begin{theorem}\label{HY-THEOREM_2}
A linear operator $\linA$ is the infinitesimal generator of a C$_0$-semigroup
{\SG} of type $(1,\omega)$ iff

\begin{itemize}[leftmargin=15pt]
\item[(a)]
$\linA$ is densely defined and closed.
\item[(b)]
(ii) The resolvent set $\rho(\linA)$ contains the interval
$\J:=(\omega,\infty)$,
and for $\zeta \xeof \J$ the estimate
$\norm{R(\linA,\zeta)} \leq (\zeta-\omega)^{-1}$ holds.
\end{itemize}
\end{theorem}

\subparagraph{}
\Proof:
If $\linA$ is an infinitesimal generator,
then the validness of conditions (a) and (b) was already quoted
in corollary 5.2.

\subparagraph{}
Reversely,
we want to sketch how the assertions (a) and (b) imply
that $\linA$ generates a quasi-contractional C$_0$-semigroup.
First of all,
it can be verified that \eqref{RESOLVENT_ESTIMATION_HY} implies
$
\lim_{\,\zeta \to \infty} \, \zeta R(\linA,\zeta) \eq x
$
for all $x \xeof \domain{\linA}$.
As a consequence,
the so-called \textit{Yosida operators}
\[
\linA_{n} := n \linA (n-\linA)^{-1}
               = n^2 (\linA-n)^{-1} - n
\]
approximate the operator $\linA$ as $n \xto \infty$ in the sense that
\[
\lim_{\,n \to \infty} \, \linA_n x \= \linA x
\]
for all $x \xeof \domain{\linA}$.
Since the operators $\linA_n$ are bounded,
one may define semigroups for each of them by letting
\[
e^{t \linA_n}\,x \eqdef
\sum_{k=0}^{\infty} \frac{(t \linA_n)^k}{k!}\,x
\quad \quad (x \xeof \banX).
\]
It follows that $\snorm{e^{t \linA_n}} \leq 1$ for $t > 0$,
so we obtain even semigroups of contractions.
Now the main part of the proof is to show that
this sequence of semigroups is Cauchy and the semigroup generated
by the operator $\linA$ is given by
\begin{equation}\label{YOSIDA_SEMIGROUP}
S(t)x \eqdef \lim_{\,n \to \infty} e^{t \linA_n}\,x
\quad \quad (x \xeof \banX).
\end{equation}
To be precise,
one first has to verify property (S-1) and (S-2)
and second the operator $\linA$ to be
the infinitesimal generator of the semigroup \eqref{YOSIDA_SEMIGROUP}.
% --tbd
\qed

\paragraph{}
We now come back to the proof of theorem \ref{GENERAL_HY-THEOREM}
by using the assertion of the already proved theorem \ref{HY-THEOREM_2}.
% --tbd
\qed

\paragraph{}
Recall that in the more special case,
where $M \eq 1$ and $\omega \eq 0$,
the semigroup of type $(1,0)$ is called a \textit{semigroup of contractions}.
From theorem \ref{HY-THEOREM_2} we deduce the following

\begin{corollary}\label{HY-THEOREM_3}
Let $\linA$ be a linear operator acting in a complex Banach space $\banX$
with domain of definition $\domain{\linA}$.
Then the following assertions are equivalent to each other:

\begin{itemize}[leftmargin=15pt]
\item[(a)]
$\linA$ generates a strongly continuous semigroup of type $(1,0)$.

\item[(b)]
$\linA$ is densely defined and closed,
the resolvent set of $\linA$ fulfills $(0,\infty) \xsubset \rho(\linA)$
and the resolvent function $R(\linA,\zeta)$ satisfies
\[
\norm{R(\linA,\zeta)}
\, \leq \, \zeta^{-1} \quad (\zeta > 0).
\]
\end{itemize}
\end{corollary}

\paragraph{}
The statement $(0,\infty) \xsubset \rho(\linA)$ in this corollary means that
one has to verify the existence and uniqueness of a solution
for each stationary equation $\linA x - \zeta x \eq y$,
where $\zeta > 0$ and $y \xeof \banX)$ is arbitrary.

\section{The Lumer-Phillips Theorem}
%------------------------------------------------------------------------------%
Preliminary,
the first theorem in this section deals with operators
having the \textit{bounded-below} property.

\subparagraph{} %- bounded below property
A linear operator $\linA$ acting in a Banach space $\banX$
is called \textit{bounded below}
if there is $\theta {\,>\,} 0$ such that
\begin{equation}\label{BOUNDED_BELOW_PROPERTY}
\norm{\linA x}_Y \, \geq \, \theta \norm{x}_X \fa x \xeof \domain{\linA}.
\end{equation}

\subparagraph{} %- bounded below theorem
The \textit{bounded-below theorem} is a fundamental assertion
in operator theory and is helpful to verify the closeness of $\range{\linA}$:

\begin{theorem}\label{BOUNDED_BELOW_THEOREM}
If $\linA$ fulfills \eqref{BOUNDED_BELOW_PROPERTY},
then $\linA$ is injective and the operator
$\linA^{-1} \: \range{\linA} \xsubset \banX \xto \banX$ is bounded
with operator norm $\norm{\linA^{-1}} \leq \theta^{-1}$.
Furthermore,
the range $\range{\linA}$ is closed \iff $\linA$ is closed.
\end{theorem}

\subparagraph{}
\Proof:
Injectivity of the operator $\linA$ directly follows
from \eqref{BOUNDED_BELOW_PROPERTY}.
To show that $\range{\linA}$ is closed,
let $y_n \to y$ and $y_n \xeof \range{\linA}$ for $n \in \N$.
Then for each $n \xeof \N$
there is $x_n \xeof \domain{\linA}$ such that $y_n \eq \linA x_n$.
In fact,
$\{x_n\}$ is a Cauchy sequence by \eqref{BOUNDED_BELOW_PROPERTY}
with limit element $x \xeof \banX$.
Since $\linA$ is closed,
we have $x \xeof \domain{\linA}$ and $\linA x \eq y$.
Hence $y \in \range{\linA}$ and $\range{\linA}$ is closed.

The inverse operator $\linA^{-1} \: \range{\linA} \xto \banX$
is closed as well,
hence by the closed graph theorem it is bounded
with operator norm $\theta^{-1}$.
\qed

\paragraph{} %- dissipative operators
For the rest of this section we switch to the Hilbert space setting
and assume $\banX \eq \hilH$ to be a complex Hilbert space.
 
\subparagraph{}
A linear operator $\linA$ acting in $\hilH$ is said to be \textit{dissipative},
if the numerical range
$\Theta(\linA):=\{\scalar{\linA x}{x}\:x \xeof \domain{\linA}, \norm{x} \eq 1\}$
fulfills
\[
\Theta(\linA) \subseteq \C^- \eqdef \{\zeta \xeof \C \: \re \zeta \leq 0\},
\]
or equivalently,
if $\re \, \scalar{\linA x}{x} \leq 0$ for all $x \xeof \domain{\linA}$.
Furthermore,
the operator $\linA$ is said to be \textit{quasi-dissipative},
if there is $\lambda_0 \xeof \R$ such that $\linA + \lambda_0$ is dissipative.

\paragraph{} %- dissipativity property
The next theorem describes dissipativity by another condition:

\begin{theorem}\label{DISSIPATIVITY_THEOREM}
$\re \, \scalar{\linA x}{x} \leq 0$ for $x \xeof \domain{\linA}$
is equivalent to
\begin{equation}\label{DISSIPATIVITY_PROPERTY}
\norm{(\zeta - \linA)x} \, \geq \, \zeta \norm{x} \quad \quad (\zeta > 0).
\end{equation}
\end{theorem}

\subparagraph{}
\Proof:
For $x \xeof \domain{\linA}$ the inequality
\[
\norm{\zeta x - \linA x}^2 - \norm{\zeta x}^2
\, \geq \,
- 2 \zeta \; \re \, \scalar{\linA x}{x}
\, \geq \, 0
\]
holds,
which implies \eqref{DISSIPATIVITY_PROPERTY}.
\qed

\paragraph{}
Property \eqref{DISSIPATIVITY_PROPERTY} implies that the operator
$\zeta-\linA$ is bounded below,
hence $\zeta-\linA$ is injective and $\range{\zeta-\linA}$ is closed.

\paragraph{} %- m-dissipative operators
A dissipative operator $\linA$ is said to be \textit{m-dissipative},
if it does not have any proper dissipative extension.
Equivalently,
a dissipative operator $\linA$ is m-dissipative,
if there is $\zeta_0 > 0$ such that $\zeta_0-\linA$ is surjective.
The property $\range{\zeta_0-\linA} \eq \banX$ is called
the \textit{range condition} for a m-disspative operator.

\paragraph{} %- m-dissipative operators are densely defined and closed
The following theorem is used
in the proof of theorem \ref{LUMER-PHILLIPS_THEOREM} below:

\begin{theorem}\label{DISSIPATIVE_OPERATOR_IS_DENSELY_DEFINED_AND_CLOSED}
If $\linA$ is a m-dissipative linear operator,
then $\linA$ is densely defined in $\hilH$ and closed.
\end{theorem}

\subparagraph{}
\Proof:
To verify that $\linA$ is densely defined,
let $y \xeof \domain{\linA} ^{\bot}$
and $x \xeof \domain{\linA} $ be a vector
such that $x - \linA x \eq y$.
Surely,
such a vector exists since $\1-\linA$ is surjective.
Then we compute
\[
0 \= \scalar{x}{y} \=
\scalar{x}{x- \linA x} \= \scalar{\linA x}{x} - \norm{x}^2 \, \leq \, 0,
\]
which implies $x \eq 0$ and $y \eq x - \linA x \eq 0$.
Hence $\domain{\linA} ^{\bot} \eq \{\0\}$ and consequently
$\closure{ \domain{\linA} } \eq \hilH$.

\subparagraph{}
Two show that $\linA$ is closed,
we have to verify that if $x_n \to x$ and $\linA x_n \to y$,
then $x \xeof \domain{\linA} $ and $\linA x \eq y$.
Clearly,
$(\zeta_0 - \linA)x_n \to \zeta_0 x-y$, hence
\[
x \= \lim_{n \to \infty} x_n \=
\lim_{n \to \infty} (\zeta_0 - \linA)^{-1}(\zeta_0 - \linA) \, x_n
\=
(\zeta_0 - \linA)^{-1}(x-y)
\]
so
$x \eq (\zeta_0 - \linA)^{-1}(x-y)
\, \Longleftrightarrow \, \zeta_0 x - \linA x \eq \zeta_0 x - y
\, \Longleftrightarrow \, \linA x \eq y$.
\qed

\paragraph{}
Our next theorem provides a very useful property for dissipative operators:

\begin{theorem}
If $\linA$ is m-dissipative,
then $\zeta - \linA$ is even surjective for all $\zeta > 0$,
hence $\rho(\linA) \eq (0,\infty)$.
\end{theorem}

\subparagraph{}
\Proof:
Consider the set
\[
\Lambda :=
\{ \zeta \xeof \R \: 0 < \zeta < \infty \and
\range{\zeta-\linA} \eq \hilH \}
\]
and notice that $\Lambda \neq \emptyset$, since $\zeta_0 \xeof \Lambda$.
Now,
if $\zeta \xeof \Lambda$, then $\zeta \xeof \rho(\linA)$
by property \eqref{DISSIPATIVITY_PROPERTY},
Since $\rho(\linA)$ is open,
any neighborhood of $\zeta$ is a subset of $\rho(\linA)$.
The intersection of this neighborhood with the real line
is clearly in $\Lambda$ and therefore $\Lambda$ is open.
On the other hand,
let $\zeta_n \xeof \Lambda$ with $\zeta_n \xto \zeta > 0$.
For every $y \xeof \hilH$ there exists an element $x_n \xeof \domain{\linA}$
such that
\begin{equation}\label{VVV}
\zeta_n x_n - \linA x_n \eq y.
\end{equation}
By using theorem \ref{DISSIPATIVITY_THEOREM},
we deduce that
$\norm{x_n} \leq \zeta_n^{-1} \norm{y} \leq c$ for some $c > 0$.
We notice now that
\[
\zeta_m \norm{x_n - x_m} \, \leq \, \norm{\zeta_m(x_n - x_m) - \linA(x_n - x_m)}
\]
\[
= \abs{\zeta_n-\zeta_m} \norm{x_n} \, \leq \, c \cdot \abs{\zeta_n-\zeta_m} .
\]
Therefore $\{x_n\}$ is a Cauchy sequence.
Let $x_n \xto x$.
Then by \eqref{VVV} $\linA x_n \xto \zeta x - y$.
Since $\linA$ is closed,
$x \xeof \domain{\linA}$ and $\zeta x - \linA x \eq y$.
Therefore $\range{\zeta - \linA} \eq \hilH$ and $\zeta \xeof \Lambda$.
Thus $\Lambda$ is closed in $(0,\infty)$.
Because $\Lambda$ is non-empty as well as open and closed at the same time,
we have $\Lambda \eq (0,\infty)$.
\qed

\paragraph{} %- estimate for the inverse of a m-dissipative operator
Finally,
for all $\zeta > 0$ and $y \xeof \range{\zeta-\linA}$
the following estimation is valid:
\[
\snorm{(\zeta-\linA)^{-1}y} \, \leq \, \zeta^{-1} \norm{y}.
\]
This is a consequence of the fact that $\zeta - \linA$ is closed
and of the bounded-below theorem,
which quotes that $\range{\zeta-\linA}$ is closed.
In fact,
in this case the inverse operator
\[
(\zeta - \linA)^{-1} : \range{\zeta-\linA} \to \domain{\linA}
\]
is bounded and has operator norm
$\norm{(\zeta-\linA)^{-1}} \leq \zeta^{-1}$,
so the operator $(\zeta-\linA)^{-1}$ forms a contraction for all $\zeta \geq 1$.

\paragraph{} %- Lumer-Phillips theorem
Let us now switch to the main theorem of this section.
By means of the notation of m-dissipativity,
the \textit{Lumer-Phillips} \textit{theorem} classifies the linear operators
being the infinitesimal generators of contractive C$_0$-semigroups.

\begin{theorem}\label{LUMER-PHILLIPS_THEOREM}
Let $\linA$ be a linear operator in a complex Hilbert space.
Then $\linA$ is m-disspipative
\iff
$\linA$ is the infinitesimal generator of a contractive C$_0$-semigroup.
\end{theorem}

\subparagraph{}
\Proof:
Let $\linA$ be m-dissipative.
Then theorem \ref{DISSIPATIVE_OPERATOR_IS_DENSELY_DEFINED_AND_CLOSED} quotes
that $\linA$ is densely defined and closed.
By theorem \ref{HY-THEOREM_3},
it is also to show that
$(0,\infty) \xsubset \rho(\linA)$ and the resolvent estimate holds,
that is,
for every $\zeta > 0$ we have $\norm{R(\linA,\zeta)} \leq \zeta^{-1}$.
In fact,
since $\range{\zeta - \linA} \eq \hilH$ for $\zeta > 0$,
inequality \eqref{DISSIPATIVITY_PROPERTY} implies
that $(\zeta-\linA)^{-1}$ is a bounded linear operator,
hence $\zeta \xeof \rho(\linA)$.
Moreover,
the resolvent estimate follows from \eqref{DISSIPATIVITY_PROPERTY}.
Thus the Hille-Yosida theorem yields that
$\linA$ is the infinitesimal generator of a C$_0$-semigroup of contractions.

\subparagraph{}
Reversely,
if $\linA$ is the infinitesimal generator of a C$_0$-semigroup of contractions,
then by theorem \ref{HY-THEOREM_3} we have $(0,\infty) \xsubset \rho(\linA)$
and therefore $\range{\zeta-\linA} \eq \hilH$ for all $\zeta > 0$.
In addition,
if $x \xeof \domain{\linA}$,
then
\[
\abs{\scalar{S(t)x}{x} } \, \leq \, \norm{S(t)x} \norm{x} \, \leq \, \norm{x}^2
\]
and
\[
\re \, \scalar{S(t)x - x}{x} \= \re \, \scalar{S(t)x}{x} - \norm{x}^2
\, \leq \, 0.
\]
Dividing this line by $t > 0$ and letting $t \downarrow 0$ yields
$\re \, \scalar{\linA x}{x} \leq 0$,
which shows that $\linA$ is dissipative.
In summary, the operator $\linA$ proves to be m-dissipative.
\qed

\paragraph{} 
Recall that to each densely defined linear operator $\linA$
acting in a Hilbert $\hilH$
there is a linear operator $\linA^*$,
called the \textit{adjoint operator} of $\linA$,
such that
$
\scalar{\linA x}{y} \eq \scalar{x}{\linA^* y}
$
for all $x \xeof \domain{\linA}$ and $y \xeof \domain{\linA^*}$.

\subparagraph{}
Having the concept of Hilbert adjoint available,
we obtain a further criterion for $\linA$ to be
the generator of a contraction C$_0$-semigroup: 

\begin{theorem}\label{ADJ_DISS}
Let $\linA$ be a densely defined linear operator
in a complex Hilbert space $\hilH$.
If both,
$\linA$ and $\linA^*$ are dissipative,
then the operator $\linA$ is the generator
of a C$_0$-semigroup of contractions.
\end{theorem}

\Proof.
From theorem \ref{LUMER-PHILLIPS_THEOREM} it is sufficient to show
that $\range{\1-\linA} \eq \banX$.
Since $\linA$ is dissipative and closed,
$\range{\1-\linA}$ is a closed subspace of $\banX$.
If $\range{\1-\linA} \neq \banX$ then there exists $x \xeof \hilH$,
$x \neq \0$ such that
$\scalar{x}{x-\linA x} \eq 0$ for $x \xeof \domain{\linA}$.
This implies $x-\linA x \eq \0$.
Since $\linA^*$ is also dissipative,
it follows from theorem \ref{DISSIPATIVITY_THEOREM}
that $x \eq \0$,
contradicting the construction of $x$.
\qed

\paragraph*{} %- the Lumer-Phillips theorem in a Banach space
We have formulated the Lumer-Phillips theorem
in the setting of a complex Hilbert space.
Let us finally mention that this theorem is also valid when $\linA$ is a
linear operator acting in a complex Banach space.
In order to quote an analogous theorem,
we have to substitute the inner product in a Hilbert space
by a similar mapping,
where \textit{dual pairs} are used.

\subparagraph{}
Let $\banX$ be a Banach space and $\banX^*$ its dual space,
that is,
the linear space of continuous linear functionals $x^* \: \banX \to \C$.
We denote the value of $x^* \xeof \banX^*$ at $x \xeof \banX$ by either
$\scalar{x^*}{x}$ or $\scalar{x}{x^*}$.
For any $x \xeof \banX$ let $F(x) \xsubset \banX^*$ be the \textit{duality set}
defined by
\[
F(x) \= \{x^* \in \banX^*: \scalar{x^*}{x} \eq \norm{x}^2 \eq \norm{x^*}^2\}.
\]
It follows from the \textit{Hahn-Banach theorem}
that $F(x) \neq \emptyset$ for $x \xeof \banX$.
Then a linear operator $\linA$ acting in $\banX$ is called \textit{dissipative}
if for every element $x \xeof \domain{\linA}$ there exists $x^* \xeof F(x)$
such that $\re \, \scalar{x^*}{x} \leq 0$. 

\subparagraph{}
If $\banX \eq \hilH$ is a Hilbert space,
however,
the set $F(x)$ consists of the singleton $\{x\}$
after identifying the element $x \xeof \hilH$
with the continuous linear functional
$f_x := \scalar{x}{\,\cdot\,} \xeof \hilH^*$.

\section{Applications to Partial Differential Operators}
%------------------------------------------------------------------------------%
The previously presented theory has many applications,
especially with respect to linear partial differential operators.

\begin{itemize}[leftmargin=12pt,itemsep=5pt]
\item[1.]
We consider the Laplacian $\Delta$ in the entire Euclidean space $\Rd$.
Then theorem \ref{LUMER-PHILLIPS_THEOREM} ensures
that the realization $\linA_{\Delta}$ in Hilbert space $\banX \eq L^2(\Rd)$
of square-integrable functions $u \: \Rd \xto \C$
and with domain of definition
\[
\domain{\linA_{\Delta}} \eqdef H^2(\Rd)
\]
and action $\linA u := \Delta u$ for $u \xeof \domain{\linA_{\Delta}}$
is the infinitesimal generator
of a strongly continuous semigroup {\GG},
denoted the \textit{Gauss-Weierstrass} \textit{semigroup}.
In fact,
it is not hard to show that the operator $\linA_{\Delta}$ is dissipative,
and m-dissipativity follows from self-adjointness of $\linA_{\Delta}$.

\subparagraph{} %- heat kernel function
It is a consequence of theorem \ref{LUMER-PHILLIPS_THEOREM}
that there is a uniquely defined solution $u \xeof C^1(\J \xtimes \Rd,\R)$
for the homogeneous Cauchy problem
\[
\begin{cases}[1.2]
\quad \dot{u} \= \linA_{\Delta} u,\\
\quad u(x,0) \= u_0(x) \for x \in \Rd\\
\end{cases}
\]
with initial mapping $u_0 \xeof H^2(\Rd)$.
Furthermore,
it is possible to represent the semigroup {\GG}
explicitly by means of the \textit{heat kernel} function
\[
\Phi_d(t,x) \eqdef \frac{1}{(4 \pi t)^{d/2}} \; e^{- \abs{x} ^2/4t},
\]
that is,
for all $t > 0$ and $x \xeof \Rd$ we have
\[
(G(t) \, u_0) \, (x)
\=
\frac{1}{(4 \pi t)^{d/2}}
\int_{\,\Rd} e^{- \frac{\abs{x-y} ^2}{4t}} u_0(y) \, dy
\=
(\Phi_d(t) \ast u_0)\,(x).
\]

\item[2.]
Next we consider the Laplacian $\Delta$ in a bounded domain $\G \xsubset \Rd$.
The differential expression gives rise to a linear operator $\linA_{\Delta}$
in Hilbert space $\Lh$ and with domain of definition
\[
\domain{\linA_{\Delta}} \eqdef
\{u \in \Lh\: \,u \in \HN{1}{\G}, \, \Delta u \in \Lh\},
\]
that is, homogeneous Dirichlet boundary conditions are incorporated into
the domain of definition of $\linA_{\Delta}$. 
If the boundary of $\G$ is sufficiently smooth,
then it follows from elliptic regularity theory that actually
$\domain{\linA_{\Delta}} \eq \HN{1}{\G} \cap H^2(\G)$ holds.

\subparagraph{}
By means of \textsc{Green}'s formula it can be shown that
the operator $\linA_{\Delta}$ is dissipative.
Moreover,
using Hilbert space methods and especially the Lax-Milgram lemma,
$\linA_{\Delta}$ fulfills the range condition
and hence proves to be m-dissipative.
Thus it is the infinitesimal generator of a C$_0$-semigroup.

\subparagraph{}
In conclusion, 
the homogeneous Cauchy problem
\[
\begin{cases}[1.2]
\quad \dot{u} \= \linA_{\Delta} u,\\
\quad u(x,t)  \= 0      \for t > 0 \and x \in \bndG,\\
\quad u(x,0)  \= u_0(x) \for x \in \clG\\
\end{cases}
\]
has a unique solution $u \xeof C^1(\J \xtimes \clG,\R)$.
But for general domains it is not possible to represent the semigroup
by a kernel function.

\item[3.]
Let $\banX = L^2(\R)$ and $\linA u:= i\,u''$
for $u \xeof \domain{\linA} :=  H^2(\R)$.
Since the operator $u \xeof H^2(\R) \mapsto u''$ is self-adjoint,
it follows that the operator $\linA$ and its adjoint $\linA^*$ fulfill
the properties
$\scalar{\linA u}{u} \xeof i\R$ and $\scalar{\linA^* u}{u} \xeof i\R$,
hence they are dissipative.
Thus the operator $\linA$ is the generator of a contraction semigroup
by theorem \ref{ADJ_DISS}.

\item[4.]
Let $\banX = L^2(0,1)$ and $\linA u := u'$
for $u$ in the domain of definition
\[
\domain{\linA} \eqdef \{u \in H^1(0,1) \: u(1) = 0\}.
\]
Then $\linA$ is m-dissipative and hence the generator of a C$_0$-semigroup
of contractions.
Indeed, we have
\[
\scalar{\linA u}{u} \= \int_0^1 u' u \, dx \= - \frac{1}{2} u^2(0) \,\leq\, 0.
\]
To verify the range condition, consider the differential equation
$u -u' = f$ for $f \xeof L^2(0,1)$ with solution
\[
u(t) \= e^t - \int_0^t f(s)\,ds, \quad (t \in [0,1]).
\]

\item[5.]
Let $\banX = C_{ub}(\R)$,
the Banach space of bounded and uniformly continuous functions $u \:\R \xto \R$
and $\linA u := u'$ for $u$ in the domain of definition
\[
\domain{\linA} \eqdef
\{u \in C_{ub}(\R) \: u \in C^1(\R) \and u' \in C_{ub}(\R)\}.
\]
Then $\linA$ is the generator of the \textit{translation} C$_0$-semigroup
{\TG} in $\R$ given by $T(t)u(s) := u(s+t)$
for $u \xeof C_{ub}(\R)$ and $t,s \xeof \R$.
\end{itemize}

%------------------------------------------------------------------------------%
%                                 BIBLIOGRAPHY                                 %
%------------------------------------------------------------------------------%
\newpage
\thispagestyle{empty}

\vspace{8mm}
\begin{thebibliography}{99}
\begin{small}

\bibitem{EN}  K.-J.\,Engel, R.\,Nagel:
\textit{A Short Course on Operator Semigroups}.
Universitext, Springer, 2006.

\bibitem{Go}  J.\,Goldstein:
\textit{Semigroups of Linear Operators and Applications}.
Oxford University Press, 1985.

\bibitem{Ka} T.\,Kato:
\textit{Perturbation Theory for Linear Operators}, Springer Verlag.
Grundlehren der mathematischen Wissenschaften, Band 132, 1995.

\bibitem{Pa}  A.\,Pazy:
\textit{Semigroups of Linear Operators
        and Applications to Partial Differential Equations}.
Applied Mathematical Sciences Vol.\,44, Springer, 1983.

\bibitem{RR}  M.\,Renardy, R.\,Rogers:
\textit{An Introduction to Partial Differential Equations}.
Texts in Applied Mathematics 13, Springer, 2004.

\bibitem{Yo}  K.\,Yosida:
\textit{Functional Analysis}.
Springer, 1980.

\end{small}
\end{thebibliography}
%------------------------------------------------------------------------------%
%                               END OF DOCUMENT                                %
%------------------------------------------------------------------------------%
\end{document}
